\section{Source audit, proposed status changes, and further research}
\label{sec:audit-research}

\subsection{What should change in the existing programme}

The audit was directed at the statements used or continued in this
article. It was not a line-by-line certification of every chapter in the
large collective manuscript. No mathematical counterexample or invalid
proof step was established in those source statements. The principal
proposed changes are strengthened theorems and more precise boundaries
between formal and arithmetic conclusions.

\begin{enumerate}[leftmargin=*,label=\arabic*.]
\item \textbf{Split the optimal-prefix conjecture by scope.}
In \cite[\texttt{conj:firstq}]{CanonicalFifth}, replace the $p=2$,
$p=4$, and $p=6$ parts by Theorems~\ref{thm:rank-two}, \ref{rf:main},
and \ref{rs:main},
including their determinants and Corollary~\ref{cor:ordinary-two}.
Retain the diagonal theorem already proved there. Theorem
\ref{thm:general-small-determinant} gives an explicit fixed-size
reduction for all even pairs. Leave unequal even pairs $8\le p<q$
as the remaining conjecture. Finite checks at those orders should continue to be identified
as checks, with no inferred all-order conclusion.

\item \textbf{Extend the directional-zeta chapter through cubic order.}
Append Theorem~\ref{tc:main} and the two explicit cancellations
\eqref{tc:two-ray-example} and \eqref{tc:example} after the inherited
quadratic theory. Preserve the precise admissible-ray condition and
the displayed polar subtraction. A global third Taylor tensor of $T$
at the origin is not defined; the theorem concerns specified one-variable
restrictions. The kernel theorem should be attached to the exact local
restrictions it classifies.

\item \textbf{Replace the unstructured higher-power task by a proved
algorithm and a smaller arithmetic task.}
Theorems~\ref{dq:normalform} and \ref{dq:parity} settle the formal
integration-by-parts classification in the chosen differential-polynomial
space. Equations~\eqref{dq:canonicalsixthformula} and
\eqref{dq:canonicalseventh} give explicit applications. What remains is
the spectral evaluation of their specified mixed moments, not the
existence of an integration-by-parts procedure.

\item \textbf{Preserve the two finite-part normalizations.}
Centered tail moments use continuation of $D_{p,q}(s)$, while singular
integrals use the constant term in the original coordinate $x$.
The proofs state why the particular ordinary subtractions recover those
values. An unannounced translation of the cutoff or omission of
$\mathcal E(T)$ could alter constants. The special binomial omission
rule in \eqref{eq:entries-foundation} must also remain explicit: restoring
a forbidden lower-index-zero term changes the matrix under discussion.

\item \textbf{Keep the priority and independence statements narrow.}
The Bernoulli Hankel determinant and differential-polynomial exactness
machinery have classical antecedents. The height-one coefficient bridge
is inherited. The new content lies in the mixed determinant, the
resulting exact cancellations, the cubic directional reduction, and the
explicitly normalized higher-power applications. Formal coordinate counts
do not prove independence of zeta values or of the remaining mixed
integrals. No historical priority assertion beyond the inspected material
is needed for these mathematical contributions.
\end{enumerate}

The manuscript's short $S_6$ identity and its revised $S_8$ candidate
remain conjectural. The older rejected $S_8$ vector is a different
statement; it must not be substituted for the current candidate. The
present results neither prove nor refute those period identities
\cite[README and editorial ledger]{Repo}.

\subsection{A proposed programme of exact research questions}

The following questions are selected for identities, reduction theorems,
and explicit primitives. They are ordered approximately by their
connection to the results above, rather than by an unproved prediction
of difficulty.

\begin{enumerate}[leftmargin=*,label=\textbf{R\arabic*.}]
\item \textbf{Prove the remaining unequal-order prefix conjecture.}
For $p=2a<2b=q$ with $a\ge4$, use the explicit $a$-square matrix
in Theorem~\ref{thm:general-small-determinant} to finish the remaining
nonvanishing problem. The shifted-Hankel part is already known to be
nonsingular, and the weighted Bernoulli reduction has degree at most
$2a-1$. The next task is to evaluate the residual determinant and find
an arithmetic obstruction to its vanishing, as in the $a=1,2,3$ proofs. The exact $(a,b)=(4,9)$ example
in Section~\ref{rs:section} rules out the naive uniform valuation $-a$;
any more general argument must accommodate that obstruction.
A proof must control
nonvanishing of the final rational factor, not just display a product of
unknown minors.

\item \textbf{Determine every reconstruction row in closed form.}
Theorem~\ref{thm:rational-row-two} gives the rational coordinate $Z_0$.
An explicit formula for every row of $M_b^{-1}$ would express each even
zeta coordinate through the same optimal moment prefix. Orthogonal
polynomial kernels followed by the rank-one inverse formula offer a
concrete route. Seek coefficient formulas with finite binomial sums and
prove their normalization directly. The result would produce systematic
ordinary harmonic-tail identities with prescribed $\zeta(2k)$ values.

\item \textbf{Extend the rational tail cancellations to shifted or
character-weighted tails.}
Replace the ordinary tails by Hurwitz tails with a rational shift, or
by tails weighted by Dirichlet characters. Determine which centering and
parity permit an exact polynomial subtraction, and which generalized
Bernoulli determinants replace $M_b$. The half-integer vanishing
$\zeta(-2r,1/2)=0$ is a special ingredient of the present proof; in a new
shift it must be replaced by an explicitly computed finite correction.

\item \textbf{Find finite polylogarithmic evaluations associated with
the boundary family.}
Equation~\eqref{eq:polylog-boundary} is an exact endpoint identity for
an explicit finite differential operator. Determine whether its numerator
admits useful integrated forms at interior algebraic arguments, with
depth reduction proved by differential identities and base values.
The test is an equality of functions or a rigorously anchored primitive,
not a numerical integer relation at a single argument.

\item \textbf{Determine the arithmetic meaning of the second cubic
Tornheim coordinate.}
The constant $\kappa$ is fixed by $J_{AA}(0)$ in the canonical
subtraction, and \eqref{tc:kappa-Euler} already expresses it through
the strict Euler double-zeta ray $\zeta(2t,t)$. Starting from that
proved representation, compare it with the inherited cubic digamma
invariant and seek a reduction. The concrete question is whether
$\kappa$ belongs to a specified algebra generated by ordinary zeta
derivatives and $\Omega$, or whether a further higher-depth coordinate
is required. A coordinate change alone does not answer that question.

\item \textbf{Compute the quartic and quintic ray cancellation laws.}
Theorem~\ref{tc:kernel} gives formal ambiguity dimensions $3$ and $6$
at degrees four and five. Construct explicit homogeneous direction
functions for these spaces, derive the ordinary-zeta inhomogeneous part,
and find short integer-weighted combinations of positive rays that
annihilate all undetermined coefficients. These are finite symbolic
problems once the analytic polar normalization has been fixed, but the
coefficient calculation must be independently checked.

\item \textbf{Add further Tornheim functional relations to the local
kernel and compute the exact reduction.}
The kernel theorem deliberately uses a specified subset of the known
relations. Linearize additional functional equations from the classical
Tornheim literature at the same origin, keeping their Gamma and trigonometric
factors to the required order. Determine whether they reduce the cubic
two-coordinate space or only change its basis. At higher degrees, compute
the quotient of the kernel by these additional restrictions before making
claims about necessary arithmetic data.

\item \textbf{Evaluate the retained sixth mixed moment in an explicitly
stated spectral family.}
The quantity $\widetilde M_6$ has a convergent subtracted integral in
\eqref{dq:M6ordinary}. Use its pole principal parts, rather than numerical
fitting, to express it through ordered harmonic zeta derivatives. Then
test whether the resulting combination reduces within a fixed family of
height-one, depth-two, or Tornheim coordinates. The desired theorem is an
explicit evaluation or a proved reduction in that family; the present
formula makes no assertion that $\widetilde M_6$ is irreducible.

\item \textbf{Determine the exact image of the differential quotient.}
The partition theorem classifies formal integration by parts. Define the
map from this quotient, together with the endpoint functional, to the
harmonic spectral coordinates supplied by the pole expansion. Compute
its kernel at each weight and identify structural generators for the
additional relations. The square formula is an example of a relation
available after substitution into the digamma function that is not an
identity in the free differential-polynomial quotient.

\item \textbf{Combine higher Stieltjes coefficients with the normal
form and parameter primitives.}
The fixed Mellin family in Proposition~\ref{prop:height-one-foundation}
gives every $\lambda_{d,j}$ and ordinary $\gamma_j$ at depth one.
Introduce the appropriate logarithmic parameter derivatives into the
Gamma family, determine their differential normal forms, and retain the
complete endpoint terms. A useful target is a closed recurrence for
primitive families whose spectral derivatives isolate one prescribed
$\lambda_{d,j}$, followed by analogous formulas for Hurwitz
$\gamma_j(a)$ with a fixed shift convention.

\item \textbf{Explain the coefficient patterns of the Gamma normal
forms.}
The Appell property already proves, for example, that the coefficient of
$X_1^{m-4}X_2^2$ is $-\binom m4$: its undifferentiated core is the
$-X_2^2$ term in $N_4$. The remaining coefficient problem is therefore
concentrated in $K_j=N_j|_{X_1=0}$. Prove general formulas for families
of repeated-largest partitions in these cores using the finite recurrence.
Distinguish allowed parity coordinates from coordinates whose coefficient
actually vanishes; an all-degree support theorem would sharpen the
current formal upper count.

\item \textbf{Return to the short Gaussian identities with a fixed
period target.}
For $S_6$ and revised $S_8$, retain the repository's exact frozen
vectors and the existing shuffle equivalences. Seek a new functional or
iterated-integral identity whose boundary evaluation is exactly that
vector. A nonzero projection in a chosen formal relation space does not
refute the numerical period identity, and a small computed residual does
not prove it. The new determinant and differential methods may suggest
useful auxiliary identities, but no direct implication is established here.
\end{enumerate}

Each question specifies the missing mathematical step. In particular,
the status of a constant evaluation should change only when an analytic
identity or an exact algebraic reduction with justified boundary values
has been supplied.
